\documentclass[12pt]{jlreq}
\usepackage{amsmath, amssymb}

\begin{document}

% \examyear{2013}
% \examgroup{B}
% \examnumber{2}
% \examfield{代数}

$A = ¥mathbf{C}[T]$ を複素数体 $¥mathbf{C}$ 上の1変数多項式環とし，$K = ¥mathbf{C}(T)$ をその分数体とする．$B$ を $¥mathbf{C}$ 上の3変数多項式環 $¥mathbf{C}[X,Y,Z]$ のイデアル $I = (XY,YZ,ZX)$ による商環 $¥mathbf{C}[X,Y,Z]/I$ とする．また $¥mathbf{C}$ 上の環準同型 $f : A ¥to B$ を $f(T) = ¥overline{X+Y+Z}$ で定める．ただしここで，$¥overline{X+Y+Z}$ は $B$ における $X+Y+Z$ の同値類である．以下の問に答えよ．

(1) $B$ を $f$ により $A$ 加群とみるとき，$B$ は自由 $A$ 加群であることを示し，その階数を求めよ．

(2) テンソル積 $B ¥otimes_A K$ のべき等元をすべて求めよ．(ただしここで，$x ¥in B ¥otimes_A K$ がべき等元であるとは，$x^2 = x$ を満たすことである．)

(3) 包含写像 $A ¥to K$ がひきおこす単射 $B ¥to B ¥otimes_A K$ により $B$ を $B ¥otimes_A K$ の部分環と考え Sax，$B ¥otimes_A K$ のすべてのべき等元で $B$ 上生成される $B ¥otimes_A K$ の部分環を $C$ で表す．$C/B$ の $¥mathbf{C}$ 上の線形空間としての次元を求めよ．

(4) 乗法群の包含写像 $B^¥times ¥to C^¥times$ が同型かどうか判定せよ．

\end{document}